\section{总结与延伸}

\subsection{本讲核心要点}

本讲是 KAIST CS492D 课程的导论，系统性地建立了生成模型的理论框架：

\begin{enumerate}
    \item \textbf{生成模型的核心问题}：给定数据样本，学习底层数据分布 $p_{\text{data}}(x)$，从而生成新样本
    \item \textbf{统一框架}：所有深度生成模型都在学习从简单分布 $p(z)$ 到数据分布 $p_{\text{data}}(x)$ 的映射
    \item \textbf{GAN}：通过生成器-判别器的对抗博弈（minimax 优化）学习映射。理论优美但训练极不稳定
    \item \textbf{VAE}：通过变分推断学习条件分布 $p_\theta(x|z)$，最大化 ELBO。训练稳定，为 Diffusion Model 奠定基础
    \item \textbf{数学工具}：逆变换采样、Box-Muller 变换、reparameterization trick、边际化、贝叶斯定理、KL 散度、Jensen 不等式
\end{enumerate}

\subsection{GAN vs. VAE 对比}

\begin{table}[H]
\centering
\begin{tabular}{lll}
\toprule
特性 & GAN & VAE \\
\midrule
训练方式 & Minimax 对抗训练 & 最大化 ELBO \\
训练稳定性 & 不稳定，容易失败 & 相对稳定 \\
生成质量 & 高（尤其是图像清晰度） & 中等（倾向模糊） \\
多样性 & 容易 mode collapse & 较好 \\
理论基础 & 博弈论 & 变分推断 \\
是否有 encoder & 无（仅 generator + discriminator） & 有（encoder + decoder） \\
latent 空间结构 & 无显式约束 & 有 KL 正则化约束 \\
与 Diffusion 关系 & 弱 & 强（ELBO 框架直接推广） \\
\bottomrule
\end{tabular}
\caption{GAN 与 VAE 的系统对比}
\end{table}

\subsection{展望：Diffusion Model 与 Flow Model}

本讲为后续课程内容铺垫了所有必要的理论基础。在接下来的讲次中，课程将深入探讨：

\begin{itemize}
    \item \textbf{Diffusion Model}：可以理解为一种``多步 VAE''，通过逐步加噪和去噪来建模数据分布。ELBO 的推导可以直接推广到 diffusion 设定下
    \item \textbf{Flow Model}：通过可逆变换直接建模精确的对数似然（不需要 ELBO 这样的下界近似）
    \item \textbf{Score-based Model}：从另一个角度——学习数据分布的 score function（对数密度的梯度）——来实现生成
\end{itemize}

本讲的采样理论（逆变换采样、reparameterization trick）、概率框架（prior / likelihood / posterior）以及 ELBO 推导将在后续每一讲中反复出现和深化。

\subsection{拓展阅读}

\begin{itemize}
    \item Goodfellow et al., \textit{Generative Adversarial Nets}, NeurIPS 2014
    \item Kingma \& Welling, \textit{Auto-Encoding Variational Bayes}, ICLR 2014
    \item Kingma \& Welling, \textit{An Introduction to Variational Autoencoders}, Foundations and Trends in ML, 2019
    \item Doersch, \textit{Tutorial on Variational Autoencoders}, arXiv:1606.05908
    \item KAIST CS492D 课程主页：\url{https://diffusion.kaist.ac.kr/}
\end{itemize}

%% --- Fin del contenido --- %%

\end{document}
